\documentclass[../../main.tex]{subfiles}% 注意这里的文件路径不能用 ./main.tex ,否则用latexmk编译子文件会报错
\graphicspath{{\subfix{./image/}}} % 指定图片目录,后续可以直接使用图片文件名
% 注意这里的文件路径不能用 ../../image/ ,否则用latexmk编译子文件会报错

% 例如:
% \begin{figure}[H]
% \centering
% \includegraphics[scale=0.3]{图.png}
% \caption{}
% \label{figure:图}
% \end{figure}
% 注意:上述\label{}一定要放在\caption{}之后,否则引用图片序号会只会显示??.

\begin{document}

\section{曲线的挠率和Frenet公式}

挠率刻画曲线偏离平面的程度，即次法向量的变化率。
\begin{definition}[挠率]
设曲线的Frenet标架为$\{\bm{\alpha},\bm{\beta},\bm{\gamma}\}$，定义**挠率**
\[
\tau(s) = -\bm{\gamma}'(s)\cdot\bm{\beta}(s).
\]
挠率的绝对值$\big|\tau(s)\big|=\big|\bm{\gamma}'(s)\big|$。
\end{definition}

\begin{theorem}
非直线曲线为**平面曲线**的充要条件是其挠率$\tau(s)\equiv 0$。
\end{theorem}


Frenet标架的运动方程，是曲线论的核心公式：
\[
\begin{cases}
\bm{\alpha}'(s) = \kappa(s)\bm{\beta}(s),\\
\bm{\beta}'(s) = -\kappa(s)\bm{\alpha}(s) + \tau(s)\bm{\gamma}(s),\\
\bm{\gamma}'(s) = -\tau(s)\bm{\beta}(s).
\end{cases}
\]


\[
\tau(t) = \dfrac{\big(\bm{r}'(t),\bm{r}''(t),\bm{r}'''(t)\big)}{\big|\bm{r}'(t)\times\bm{r}''(t)\big|^2}.
\]
其中$(\cdot,\cdot,\cdot)$为向量的混合积。

\begin{example}
圆螺旋线$\bm{r}(t)=(a\cos t,a\sin t,bt)$的挠率为常数：
\[
\tau = \dfrac{b}{a^2+b^2}.
\]
\end{example}

\begin{exercise}
\begin{enumerate}
\item 计算§2.3习题中各曲线的挠率。
\item 证明：若曲线的密切平面恒过定点，则曲线为平面曲线。
\end{enumerate}
\end{exercise}




\end{document}